\documentclass[11pt]{article}

\usepackage[margin=1in]{geometry}
\usepackage{amsmath,amssymb,amsfonts,mathtools}
\usepackage{booktabs,multirow}
\usepackage{enumitem}
\usepackage{xcolor,graphicx}
\usepackage{hyperref,microtype}

\newcommand{\E}{\mathbb{E}}
\newcommand{\Var}{\operatorname{Var}}
\newcommand{\Bern}{\operatorname{Bernoulli}}
\newcommand{\sign}{\operatorname{sign}}
\newcommand{\clip}{\operatorname{clip}}
\newcommand{\MSE}{\operatorname{MSE}}
\newcommand{\KL}{D_{\mathrm{KL}}}
\newcommand{\Acc}{\operatorname{Acc}}
\newcommand{\pbit}{\text{p-bit}}
\newcommand{\pbits}{\text{p-bits}}
\newcommand{\ours}{\textsc{SPW}}

\title{Single-P-Bit Weight Encoding:\\Implementation and Experimental Execution Plan}
\author{}
\date{}

\begin{document}
\maketitle
\tableofcontents

\section{Objective and central hypothesis}
A continuous machine-learning parameter may be represented by one physical
probabilistic bit (p-bit) through its stationary probability, instead of by several
spatial binary elements. For parameter $i$ in range group $g(i)$, define
\begin{equation}
 w_i=c_{g(i)}+r_{g(i)}(2p_i-1),\qquad p_i\in[0,1].
 \label{eq:weightmap}
\end{equation}
The stochastic state is $z_i^{(t)}\sim\Bern(p_i)$, or
$s_i^{(t)}=2z_i^{(t)}-1\in\{-1,+1\}$. From $T$ observations,
\begin{equation}
 \hat w_i^{(T)}=c_{g(i)}+r_{g(i)}\frac1T\sum_{t=1}^T s_i^{(t)}.
 \label{eq:temporalestimate}
\end{equation}
The question is whether this spatial-to-temporal precision exchange can compete with
multi-p-bit encodings while preserving task quality under realistic cost, latency, and
device nonidealities. Probability as a scalar code is established; the research target
is whether one programmable probabilistic element per parameter is useful for neural
computation.

\section{Terminology and research questions}
We call the representation Single-P-Bit Weight Encoding (\ours). Keep three execution
regimes distinct: (R1) repeatedly sample, reconstruct weights, then use conventional
operators; (R2) sample stochastic weights and execute operators directly; (R3) reformulate
operators to accumulate useful expectations without explicit weight reconstruction.

\begin{description}[leftmargin=!,labelwidth=2.3cm]
\item[RQ1 Representation] How does scalar reconstruction error vary with sample budget?
\item[RQ2 Localization] How do layer, channel, block, or group ranges reduce samples?
\item[RQ3 Robustness] How much stochastic error can each network tolerate?
\item[RQ4 Tradeoff] When does one p-bit over time beat $K$ spatial p-bits?
\item[RQ5 Execution] Can useful neural computation avoid deterministic weight recovery?
\item[RQ6 Training] Can probability parameters be optimized directly?
\item[RQ7 Hardware] Do sampling, programming, calibration, and correlations preserve an advantage?
\end{description}

\section{Representation baselines and common interface}
First reproduce FP reference, INT4, and a $K$-bit bipolar multi-p-bit expansion
\begin{equation}
 w\approx \alpha\sum_{k=1}^{K}a_ks_k,\qquad s_k\in\{-1,+1\},\quad K\in\{2,4,8\},
\end{equation}
with $K=4$ as the primary comparison. Every encoding implements the same operations:
\texttt{encode(weight)}, \texttt{sample(encoded,n)}, \texttt{decode(samples)}, and
\texttt{forward(input,encoded,n)}. Do not start neural comparisons before interface parity.
Keep the implementation modular under \texttt{encodings/}, \texttt{stochastic/},
\texttt{operators/}, \texttt{training/}, \texttt{models/}, \texttt{experiments/},
and \texttt{analysis/}.

\section{Scalar simulator and precision limits}
For $w\in[-1,1]$, set $p=(w+1)/2$, draw $z_t\sim\Bern(p)$, and estimate
$\hat w_T=2T^{-1}\sum_tz_t-1$. Then
\begin{equation}
 \Var(\hat p_T)=\frac{p(1-p)}T,\qquad
 \Var(\hat w_T)=\frac{4p(1-p)}T,
 \quad \operatorname{Std}(\hat w_T)|_{p=1/2}=T^{-1/2}.
\end{equation}
Sweep $T\in\{1,2,4,8,16,32,64,128,256,512,1024,4096\}$, at least 1000 uniformly
spaced target values, and at least 1000 Monte Carlo trials per condition. Report bias,
MAE, and RMSE. Compare empirical variance to the analytical value. Define worst-case
$\sigma_T=T^{-1/2}$ and the rough effective precision
\begin{equation}
 b_{\mathrm{eff}}=-\log_2\sigma_T\simeq\tfrac12\log_2T,
 \qquad T\sim2^{2b_{\mathrm{eff}}}.
\end{equation}
This exponential temporal cost motivates direct stochastic execution rather than exact
weight reconstruction.

\section{Local ranges, robust ranges, and metadata}
For each group $G$, use $c_G=(w_G^{\max}+w_G^{\min})/2$ and
$r_G=(w_G^{\max}-w_G^{\min})/2$, then
$p_i=\tfrac12((w_i-c_G)/r_G+1)$, clipped only for numerical protection. Test
$\epsilon\in\{10^{-6},10^{-4},10^{-3},10^{-2}\}$. Grouping strategies: global,
per-layer, per-output-channel, per-input-channel, and flattened blocks of 16, 32, 64,
and 128 weights. Always count storage for every $c_G,r_G$.

Min/max ranges are outlier-sensitive. Also evaluate percentile bounds $l_G=Q_\alpha(W_G)$,
$u_G=Q_{1-\alpha}(W_G)$ at $\alpha\in\{0,10^{-4},10^{-3},5\cdot10^{-3},10^{-2}\}$.
Either clip outliers or store sparse residuals $w_i=\tilde w_i+\delta_i$; report the
residual fraction and metadata cost.

\section{Nonideal p-bits}
Support temporal Markov correlation with approximately
$\operatorname{Corr}(s_t,s_{t+\tau})=\rho^\tau$ for
$\rho\in\{0,.01,.05,.1,.25,.5,.75\}$ and compare with
$T_{\rm eff}\simeq T(1-\rho)/(1+\rho)$. Model probability calibration as
\begin{equation}
 \tilde p_i=\sigma(a_i\operatorname{logit}(p_i)+b_i),\quad
 a_i\sim\mathcal N(1,\sigma_a^2),\quad b_i\sim\mathcal N(0,\sigma_b^2),
\end{equation}
and drift as $b_i(t+1)=b_i(t)+\eta_t$, $\eta_t\sim\mathcal N(0,\sigma_{\rm drift}^2)$.
Report recalibration frequency. The hardware-equivalent simulator additionally models
finite switching rate, cross-device correlation, quantized control voltage, and readout
noise. Quantize $\hat V=Q_{B_V}(V)$ for $B_V\in\{4,6,8,10,12\}$ with
$p=\sigma(\beta V)$.

\section{Dot products and matrix multiplication}
For $y=\mathbf w^\top\mathbf x$, directly sample
\begin{equation}
 \hat y_t=\sum_i(c_i+r_is_i^{(t)})x_i,\qquad
 \bar y_T=T^{-1}\sum_t\hat y_t.
\end{equation}
Then $\E[\bar y_T]=\mathbf w^\top\mathbf x$, and for independent p-bits
\begin{equation}
 \Var(\hat y)=\sum_i x_i^2r_i^2\bigl(1-(2p_i-1)^2\bigr),\quad
 \Var(\bar y_T)=\Var(\hat y)/T.
\end{equation}
Test dimensions $d\in\{16,64,256,1024,4096,16384\}$ with Gaussian,
unit-normalized, heavy-tailed, and real activation inputs. Extend to $Y=XW$ in both
reconstructed-matrix mode ($X\hat W_T$) and sample-and-accumulate mode
($T^{-1}\sum_tXW^{(t)}$).

\section{Neural evaluation, sensitivity, and adaptive inference}
After operator gates pass, evaluate frozen FP32 MNIST MLP, Fashion-MNIST MLP/CNN, and
CIFAR-10 small CNN. Compare FP32, FP16/BF16, INT8/4, multi-p-bit $K=4$, and SPW with
global, layer, channel, and block ranges. Sweep $T\in\{1,2,4,8,16,32,64,128,256\}$.
Replace one layer at a time to estimate
$S_\ell(T)=\mathcal L(W_\ell\to\hat W_{\ell,T})-\mathcal L_{FP}$, then allocate
nonuniform $T_\ell$ according to sensitivity.

Implement empirical-standard-error early stopping at
$\tau\in\{.1,.05,.01,.005,.001\}$ and classification margin stopping. If logits
$y_{(1)},y_{(2)}$ have estimated standard deviations $\sigma_{(1)},\sigma_{(2)}$,
stop when $y_{(1)}-y_{(2)}>\gamma\sqrt{\sigma_{(1)}^2+\sigma_{(2)}^2}$.

\section{Probability-aware and Monte Carlo training}
Parameterize $p_i=\sigma(\theta_i)$ and
$w_i(\theta_i)=c_g+r_g(2\sigma(\theta_i)-1)$. Optimize task loss plus saturation
and variance regularization, e.g.
\begin{equation}
 \mathcal L_{\rm total}=\mathcal L_{\rm task}+\lambda_{\rm sat}\mathcal L_{\rm sat}
 +\lambda_{\rm var}\mathcal L_{\rm var},\qquad
 v_i=4r_g^2p_i(1-p_i).
\end{equation}
Compare expected-weight training, straight-through, Gumbel/Concrete, REINFORCE, and
local-reparameterization approaches; start with straight-through. Fine-tune only after
frozen inference is characterized.

\section{Transformers and selective conversion}
Only after small models pass, progress through a small DistilBERT-like model, GPT-2 Small,
GPT-2 Medium, a sub-billion Qwen-family model, then approximately 1B parameters. Begin
with post-training conversion. Evaluate embedding, attention Q/K/V/output, MLP up/gate/down,
and LM head separately; initially keep LayerNorm and biases deterministic.
Test hybrid layer assignments $q_\ell\in\{0,1\}$ under converted fractions
$\rho\in\{.25,.50,.75,.90,1.00\}$.

\section{Fair resource matching and Pareto results}
For $N$ parameters, a $K$-p-bit method with $T_K$ samples costs
$C_{\rm multi}=KNT_K$ stochastic-device evaluations; SPW costs $C_{\rm single}=NT_1$.
Require matched transition counts, so for $K=4$, $T_1=4T_4$. Also compare equal
physical device count, samples, estimated energy, latency, and task quality. The main
results are Pareto curves for quality against p-bits/parameter, samples, transitions,
latency, and energy, plus effective precision versus sampling and sampling versus group
granularity.

\section{Hardware cost and device counts}
Separate programming, stochastic transitions, readout, accumulation, and recalibration.
With transition energy $E_s$, bias-programming energy $E_b$, sampling frequency $f_s$,
$N_p$ devices, reuse $N_{\rm reuse}$, and sample count $T$,
\begin{equation}
 E_{\rm infer}\simeq N_pTE_s+\frac{N_pE_b}{N_{\rm reuse}},\qquad
 L\simeq T/f_s+L_{\rm program}+L_{\rm readout}.
\end{equation}
Static pretrained weights amortize bias programming across inferences. A four-p-bit
baseline uses $4N$ devices; SPW uses $N+N_{\rm metadata}$, with metadata including
range parameters and any residual table. Do not call device reduction a systems gain
unless transition, latency, metadata, and energy costs also support it.

Investigate shared range parameters per matrix, row, column, and $16\times16$,
$32\times32$, or $64\times64$ tile. If one p-bit is insufficient, use residual
multi-scale encoding $w_i=c_g+\sum_{k=1}^{K_i}r_{g,k}(2p_{i,k}-1)$ with adaptive
$K_i\in\{1,2,4\}$. Jointly optimize group-specific spatial precision
$K_g\in\{1,2,4\}$ and temporal budget $T_g\in\{1,2,4,\ldots,256\}$ under a task-loss
constraint.

\section{Information theory and end-to-end reporting}
For $T$ Bernoulli observations, Fisher information is
$\mathcal I(p)=T/[p(1-p)]$, so the Cram\'er--Rao bound gives
$\Var(\hat p)\ge p(1-p)/T$. Interpret this as the temporal precision limit and motivate
native stochastic computation and model-level tolerance, not merely replacement of an
FP16 memory cell.

Every major model reports task accuracy/perplexity, p-bits/parameter, samples, transitions,
effective precision, metadata, latency and energy estimates, and sensitivity to
calibration, correlation, and drift. Use at least five random seeds (prefer ten for
central claims), report mean and 95\% interval $\bar x\pm1.96s/\sqrt n$, and pair
evaluations on identical examples and seeds where possible.

\section{Mandatory controls and ablations}
Negative controls: variance-matched additive Gaussian weight noise; INT2/4/8; equal
numbers of deterministic and probabilistic binary elements where meaningful; and equal
total random samples. Required ablations: samples 1--4096, range grouping and block
size, min-max/percentile/learned ranges, frozen/fine-tuned/stochastic training,
correlation 0--.75, voltage precision 4--12 bits, 1/2/4 p-bits per weight,
fixed/adaptive sampling, and reconstructed/direct execution.

\begin{table}[h]
\centering
\caption{Ablation matrix}
\begin{tabular}{ll}\toprule
Factor & Values \\\midrule
Samples & $1$--$4096$ \\
Grouping / block size & global, layer, channel, block / $16,32,64,128$ \\
Range / training & min-max, percentile, learned / frozen, fine-tuned, stochastic \\
Nonidealities & correlation $0$--$.75$, voltage $4$--$12$ bits \\
Spatial / temporal / execution & $1,2,4$ p-bits; fixed/adaptive; reconstructed/direct \\\bottomrule
\end{tabular}
\end{table}

\section{Decision gates}
\begin{enumerate}[label=Gate \Alph*.,leftmargin=*]
\item Scalar Monte Carlo error agrees with theory within statistical tolerance; failure is a simulator bug.
\item Dot-product accumulation needs materially fewer effective samples than per-weight reconstruction.
\item Small-network SPW inference is within roughly 1--2 percentage points absolute accuracy at plausible sampling.
\item A pretrained Transformer retains useful quality with modest, adaptive, or fine-tuned sampling.
\item Request physical X0/XTR-0 access only after the simulator finds a favorable quality/device/transition regime.
\end{enumerate}

\section{Physical validation}
When hardware is available, characterize devices before large models. H1: calibrate requested
$p\in\{.01,.05,.1,\ldots,.9,.95,.99\}$ with up to $10^6$ observations per point. H2:
measure autocorrelation and effective sampling rate. H3: measure cross-device correlations.
H4: measure probability drift over minutes, hours, and thermal conditions. H5: replicate
synthetic dot products. H6: run one complete neural linear layer before an end-to-end network.

\section{Execution order and schedule}
The implementation order is: (1) reproduce FP/INT4/multi-p-bit baselines; (2) scalar simulator
and variance check; (3) local ranges and metadata; (4) nonideal simulator; (5) dot products;
(6) matrix multiply; (7) MNIST/Fashion-MNIST; (8) CIFAR-10; (9) sensitivity; (10) adaptive
sampling; (11) probability fine-tuning; (12) Transformer linear layers and small Transformer;
(13) spatial-temporal Pareto and cost model; (14) hardware-equivalent simulator; (15) device
characterization; (16) physical operator and network. Decision gates classify evidence and
claims; a negative result does not cancel independent experiments in the roadmap. Defer a
dependent stage only when its prerequisite is technically invalid or its required resource is
unavailable, and record that specific limitation.

Suggested experiments are \texttt{e00\_reproduce\_baseline.py}, \texttt{e01\_scalar.py},
\texttt{e02\_local\_range.py}, \texttt{e03\_nonidealities.py}, \texttt{e04\_dot\_product.py},
\texttt{e05\_matmul.py}, \texttt{e06\_mnist.py}, \texttt{e07\_cifar.py},
\texttt{e08\_adaptive.py}, \texttt{e09\_finetune.py}, \texttt{e10\_transformer.py},
\texttt{e11\_cost\_model.py}, and \texttt{e12\_hardware.py}. Preserve configs and raw
results separately from analysis and plots.

\section{Minimum serious study and failure modes}
At minimum, report the theoretical variance/information limit, scalar reconstruction, direct
dot products, one small image model and one Transformer, four-versus-one p-bit, local ranges,
sampling budgets, nonideal/correlated simulation, device/transition/latency estimates, and
matched-resource Pareto curves. The simple method is unsuccessful if useful inference needs
extreme $T$, correlation destroys precision, control-voltage precision or metadata dominates,
accuracy requires near-deterministic reconstruction, direct execution costs too much, or
recalibration is too frequent. Such results are valid findings.

\section{Potential result and paper narrative}
A strong result shows at matched quality that $N_{\pbit}^{\ours}\ll N_{\pbit}^{\rm multi}$,
while inference energy and latency remain acceptable. The paper should establish that
probability encodes the parameter, derive sampling limits, show the cost of individual
reconstruction, test whether neural aggregation reduces samples, add local ranges and
probability-aware fine-tuning, and validate nonidealities on hardware. The central design
space is
\begin{equation}
 \boxed{\text{spatial p-bits}\times\text{temporal samples}\times
 \text{local dynamic range}\times\text{network tolerance}}.
\end{equation}
One p-bit per weight is its spatially compressed extreme, not an assumed win.

\end{document}
